\begin{description}
\item[(a)] $\phi_B = \alpha = 6^\circ 41'$
  \hfill(\textit{conceptually difficult:} 5 points)

% FIGURE NEEDED: diagram of Earth and Moon showing the libration-in-latitude
% geometry (tilted Moon connected to Earth by a line through the centres);
% 2023_th_solutions.pdf page 15 (printed page 15).

\item[(b)] The maximum angle of libration $\phi_L$ occurs when the Moon is at
  point $M$ and for $e = 0.064$:

% FIGURE NEEDED: diagram of the Moon's elliptical orbit showing the Moon at the
% end of the minor axis, the second focus F2, the Earth at the first focus, and
% the distances a, b, c; 2023_th_solutions.pdf page 15 (printed page 15).

  $\phi_L$ is the Earth--Moon--F2 angle; $|EF2| = 2c$.\\
  Therefore
  $\sin(\phi_L/2) = c/a = e \implies
   \phi_L = 2\arcsin(0.064) = 7.34^\circ = 7^\circ 20'$
  \hfill(5 points)

\item[(c)] In the less accurate approximation,
  \[
    (2\phi_B + 2\phi_L)/180^\circ = x/50\% \implies x = 8\%
  \]
  and the total area visible is $S \approx 50\% + x = 58\%$.\\
  This is less accurate as it counts the overlapping areas made visible by both
  librations twice.

  A slightly more accurate approximation is given by taking the average of
  $\phi_B$ and $\phi_L$, $7^\circ 01'$. The belt made visible by libration will
  thus be $2\pi R_{\mathrm{Moon}} \cdot 7^\circ/360^\circ \approx 213$ km wide.
  The surface area of the belt is then approximately
  $213\ \mathrm{km} \times 2\pi R_{\mathrm{Moon}}
   = 2.32 \times 10^{6}\ \mathrm{km}^2$, which accounts for about 6\% of the
  lunar surface. If the students does this algebraically before substituting,
  they also get: % sic: "the students does"
  \[
    \frac{2\pi R_{\mathrm{Moon}} \cdot 2\pi R_{\mathrm{Moon}}}
         {4\pi \, 2 R_{\mathrm{Moon}}^2}
    \cdot \frac{7^\circ}{360^\circ}
    = \pi \, \frac{7^\circ}{360^\circ} \approx 6\%.
  \]
  In either case the total area visible is $S \approx 50\% + 6\% = 56\%$.
  \hfill(less accurate solution 4 points, better solution 5 points)

\item[(d)] It is necessary to find a common multiple of the period of libration
  in latitude $P_B = 27.2122$\,d (the draconic month) and the period of
  libration in longitude $P_L = 27.5545$\,d (the anomalistic month).

  The simplest way is analogous to synodic periods:
  \[
    \frac{1}{T_{\mathrm{syn}}} = \frac{1}{T_1} - \frac{1}{T_2}
  \]
  thus
  \[
    \frac{1}{T} = \frac{1}{P_B} - \frac{1}{P_L}
    \implies T \approx 2190.5\,\mathrm{d} = 74\,\mathrm{lunations}.
  \]
  Or using continued fractions:
  \[
    \frac{P_L}{P_B} = \frac{275545}{272122}
    \approx 1 + \cfrac{1}{79 + \cfrac{1}{2 + \cfrac{1}{131 + \cfrac{1}{6 + \cdots}}}}\,.
  \]
  With the first term we approximate $P_L/P_B \approx 80/79$; since
  $80 \times 27.2122 \approx 2176.9$ d and $79 \times 27.5545 \approx 2176.8$ d
  this is already close enough, and $T$ is again equal to about 74 lunations.
  \hfill(3 points)

  Since we only need half the cycle to see everything once (as the Moon swings
  away and back again during the whole cycle) therefore
  $T/2 = 37\ \mathrm{months} \approx 3$ years is enough to see all of the
  potentially visible parts of the Moon's surface.
  \hfill(2 points)
\end{description}
